\begin{figure}[H]
\centering
\begin{tikzpicture}[
 >=Stealth,
 font=\small,
 partbox/.style={
  draw=black,
  fill=white,
  rounded corners=3pt,
  line width=.65pt,
  text width=4.9cm,
  minimum height=2.35cm,
  align=center,
  inner sep=7pt
 },
 route/.style={->,line width=.85pt,draw=black}
]
\node[partbox] (measures) at (-2.95,0) {
 \textbf{Part~\ref{part:measures}}\\
 Continuous invariant\\
 full-support probabilities\\
 \Cref{thm:invariant-measures}
};
\node[partbox] (models) at (2.95,0) {
 \textbf{Part~\ref{part:spectral}}\\
 Finite-dimensional local models\\
 $\Isom(Y,e)\to\OrthogonalGroup(n)$\\
 \Cref{thm:local-models}
};
\node[partbox,minimum height=2.95cm] (ar) at (2.95,-3.55) {
 \textbf{Part~\ref{part:topology}}\\
 $\ActingGroup=\prod_{i\in\NonnegativeIntegers}\OrthogonalGroup(n_i)$\\
 $\CompactHyperspace(\BanachAmbientSpace)/\ActingGroup$ is an AR\\
 Small tracks $\Longrightarrow$ ANR\\
 Scaling contraction $\Longrightarrow$ AR\\
 \Cref{thm:part-iii-main}
};
\node[
 partbox,
 minimum height=2.95cm
] (discrete) at (-2.95,-3.55) {
 \textbf{Part~\ref{part:hilbert}}\\
 Products with finite equilateral spaces\\
 Discrete approximation\\
 \Cref{thm:discrete-approximation}
};
\node[
 partbox,
 text width=10.8cm,
 minimum height=1.9cm
] (hilbert) at (0,-7) {
 \textbf{Toru\'nczyk's characterization}\\
 AR and discrete approximation,
 with completeness and separability\\
 $\GHSpace$ is homeomorphic to real $\SeparableHilbertSpace$\\
 \textbf{Part~\ref{part:hilbert}.
 \Cref{thm:part-iv-main}}
};
\draw[route] (measures.east)--(models.west);
\draw[route] (models.south)--(ar.north)
 node[midway,left=5pt,font=\footnotesize,align=right]
 {Partition of unity\\in the Banach space $\BanachAmbientSpace$};
\draw[route] (ar.south)--(2.95,-5.5)--(0,-5.5)--(hilbert.north);
\draw[line width=.85pt,draw=black]
 (discrete.south)--(-2.95,-5.5)--(0,-5.5);
\fill[black] (0,-5.5) circle (1.5pt);
\end{tikzpicture}
\caption{The two arguments used to identify the topology of
$\GHSpace$.
Arrows indicate dependence of the proofs.
In Parts~\ref{part:measures}--\ref{part:topology},
we prove the absolute retract property.
The discrete approximation theorem in Part~\ref{part:hilbert}
uses the common preliminaries and its own metric constructions.
These two conclusions,
together with completeness and separability,
imply the Hilbert space classification.}
\label{fig:proof-strategy}
\end{figure}
